% ============================================================
% A04 — Kategori dan Konfigurasi Navbar
% ============================================================
\flowmeta{A04}{Kategori dan Konfigurasi Navbar}{Admin}

\begin{flowsection}{Tujuan}
Kelola taksonomi kategori produk yang tampil pada navbar storefront.
\end{flowsection}

\begin{flowsection}{Prasyarat}
Login sebagai admin dengan peran yang memiliki akses modul ini.
\end{flowsection}

\begin{flowsection}{Langkah}
\begin{enumerate}
  \item Buka Master Kategori (/admin/categories).
  \item Tambah kategori: nama, slug, deskripsi, urutan, tampilkan di navbar.
  \item Ubah/aktifkan kategori dari menu baris.
\end{enumerate}
\end{flowsection}

\shot{gambar/admin/a04-1.png}{Halaman Kategori dan Konfigurasi Navbar pada Panel Admin}{fig:a04-1}

\begin{flowsection}{Hasil yang Diharapkan}
Kategori tersimpan dan muncul pada navbar storefront.
\end{flowsection}

\begin{flowsection}{Penanganan Error dan Kasus Khusus}
\begin{itemize}
  \item Slug harus unik.
\end{itemize}
\end{flowsection}

\begin{flowsection}{Kaitan Modul}
\begin{itemize}
  \item \texttt{CategoryListPage.jsx}
  \item \texttt{CategoryController}
  \item \texttt{/api/categories}
\end{itemize}
\end{flowsection}

\docver{v1.0}{Sep 2026}
